% !TeX root = ../main.tex
% ---------------------------------------------------------------------------
% Soft Ground Behaviour and Face Stability
% Scope extension: not derived from the RSE3010 weekly material.
% Key sources: Broms & Bennermark (1967); Davis et al. (1980); Anagnostou &
% Kovari (1994, 1996); Mair & Taylor (1997); DAUB recommendations.
% ---------------------------------------------------------------------------
\chapter{Soft Ground Behaviour and Face Stability}
\label{ch:soft-ground}


\section{Soft ground behaviour}
\label{sec:soft-ground-behaviour}

\subsection{Soft ground as a distinct design problem}
\label{sec:soft-ground-distinct}
\scopenote{Where rock mass classification stops working. Behaviour classes:
firm, ravelling, running, flowing, squeezing and swelling ground, after
Terzaghi and Heuer.}

\subsection{Drained and undrained response}
\label{sec:drained-undrained}
\scopenote{Total and effective stress analysis in tunnelling, the role of
permeability and advance rate, and how to judge which condition governs the
short and long term.}

\section{Face stability analysis}
\label{sec:face-stability-analysis}

\subsection{Stability ratio and the classical solutions}
\label{sec:stability-ratio}
\scopenote{Broms \& Bennermark's $N$, the Davis et al.\ upper and lower bound
solutions for cohesive ground, and the cover-to-diameter dependence.}

\subsection{Wedge--silo models for face stability}
\label{sec:wedge-silo}
\scopenote{Horn's wedge with Janssen silo loading; the Anagnostou \& Kovari
formulation for frictional ground; the influence of $c'$, $\phi'$ and
groundwater on required support pressure.}

\section{Face support pressure}
\label{sec:face-support-pressure-group}

\subsection{Face support pressure design}
\label{sec:face-support-pressure}
\scopenote{Minimum pressure against collapse, maximum against blow-out and
heave, and the operating window between them. Effective versus total pressure,
and the pressure gradient over the face height.}

\subsection{Slurry support and filter cake formation}
\label{sec:slurry-support}
\scopenote{How slurry transfers pressure to the skeleton: membrane model
versus penetration model, slurry rheology, and the effect of grain size on
filter cake formation.}

\subsection{EPB support and soil conditioning}
\label{sec:epb-conditioning}
\scopenote{Chamber pressure control, plastic paste requirements, foam, polymer
and bentonite conditioning, and the applicability envelope in grading terms.}

\section{Ground loss and mixed face conditions}
\label{sec:ground-loss-and-mixed-face-conditions}

\subsection{Mechanisms of ground loss}
\label{sec:ground-loss-mechanisms}
\scopenote{Face loss, shield loss (overcut, ploughing, steering) and tail loss
(annulus, grouting, lining deflection) as separable contributions to the total
volume loss used in Chapter~\ref{ch:numerical-modelling}.}

\subsection{Mixed face conditions}
\label{sec:mixed-face}
\scopenote{Rock and soil in the same face: differential wear, steering
difficulty, over-excavation of the soft horizon and elevated settlement risk.}

\section{Alternatives to face pressure}
\label{sec:alternatives-to-face-pressure}

\subsection{Compressed air working}
\label{sec:compressed-air}
\scopenote{Principle, air loss, decompression requirements and the health and
regulatory constraints that have narrowed its use.}

\subsection{Ground improvement as an alternative to face pressure}
\label{sec:soft-ground-improvement}
\scopenote{When treating the ground is preferable to supporting the face.
Cross-reference to Chapter~\ref{ch:ground-improvement}.}
